\documentclass[10pt,a4paper]{amsart}
\usepackage[margin=19mm]{geometry}
\usepackage{amsmath,amssymb,mathtools}
\usepackage[scr=rsfs,cal=euler]{mathalfa}
\usepackage{xfrac}
\usepackage[unicode,hidelinks,bookmarks=false]{hyperref}
\newcommand{\ie}{i.e.\ }
\newcommand{\cf}{cf.\ }
\newcommand{\resp}{respectively\ }
\newcommand{\Subsection}{\subsection}
\providecommand{\nobreakdash}{\nobreak\hbox{-}}
\setlength{\emergencystretch}{2em}
\multlinegap0pt
\usepackage{bm}
\usepackage{amssymb}
\usepackage{tikz}
\newtheorem{theorem}{Theorem}
\newtheorem{lemma}{Lemma}
\newcommand{\boldx}{\boldsymbol{x}}
\title{Analysis of a finite element method for second order uniformly elliptic PDEs in non-divergence form}
\author{Weifeng Qiu}
\date{Source: arXiv:2512.14219v4 (CC BY 4.0)}
\begin{document}
\maketitle

\noindent\emph{Source and licence.} Analysis of a finite element method for second order uniformly elliptic PDEs in non-divergence form. Version arXiv:2512.14219v4 (CC BY 4.0), distributed by arXiv under the Creative Commons Attribution 4.0 International licence, http://creativecommons.org/licenses/by/4.0/, as recorded in the arXiv licence metadata for this exact version and shown on the abstract page https://arxiv.org/abs/2512.14219v4. Full licence notice: \texttt{licenses/arxiv-2512.14219-CC-BY-4.0.txt}.\par\smallskip

\noindent\textit{Editorial scope.} Counted unit: the article's theorem thm\_A\_continuous\_complete (source0.tex lines 1975--2015). The article's complete original proof of it is delivered with it (source0.tex lines 2017--2186, one \texttt{proof} environment of 170 lines). No mathematics is written, repaired or re-derived: the statement and the proof are the article's own lines, in the article's own notation, and no retained line is substituted or re-flowed.  Retained uncounted as the source-local context the proof needs: lines 139--143; lines 271--282; lines 511--515; lines 561--566; lines 569--574; lines 580--588; lines 1010--1025; lines 1695--1718; lines 1824--1840.  Nothing was omitted from any retained span, so the package carries no omission record.  No retained line contains a citation, so the wrapper carries no bibliography and no bibliography entry.  No retained line contains a figure environment, an image-inclusion command or drawing code, so no image is delivered and the package declares no asset dependency.  Editorial formatting only, in addition: an AMS \texttt{amsart} preamble that restates the article's own declarations and macros used by the retained lines, namely the environments \texttt{theorem}, \texttt{lemma} on separate counters, together with the standard packages those lines need.  The retained environments print in this document as Theorem 1, Lemma 1 in order of appearance, whereas the article prints the same items with the numbering of its own full document.  The complete native source of the article as distributed in the arXiv e-print for this exact version is bundled unchanged as \texttt{sources/191083/source0.tex}.
\begin{align}
\label{nondiv_pde_original}
A:D^{2}u + \boldsymbol{b}\cdot\nabla u + cu = f \text{ in } \Omega, \qquad 
u = 0 \text{ on } \partial\Omega.
\end{align}

\begin{align}
\label{p_range1}
& \text{(PDE analysis) } p \in (1,2]  \text{ if } \Omega  \text{ is a convex polyhedra }(d=2,3), \\
\nonumber
& \qquad \text{ while } p \in (1, \frac{4}{3} + \epsilon_{2}) \text{ if } \Omega 
\text{ is a non-convex polygon in } \mathbb{R}^{2}; \\
\label{p_range2}
& \text{(Numerical analysis) } p \in (1,2]  \text{ if } \Omega  \text{ is a convex polyhedra }(d=2,3), \\
\nonumber 
& \qquad \text{ while } p \in (\frac{4}{3}-\epsilon_{1}, \frac{4}{3} + \epsilon_{2}) \text{ if } \Omega 
\text{ is a non-convex polygon in } \mathbb{R}^{2}.
\end{align}

\begin{align}
\label{discrete_zero_order_norm}
\Vert w\Vert_{L_{h}^{p}(D)}:= \sup_{v_{h} \in V_{h}(D)} \dfrac{(w, v_{h})_{D}}{\Vert v_{h}
\Vert_{L^{p^{\prime}}(D)}},
\end{align} 

\begin{align}
\label{def_gamma}
\gamma(\boldx) = \dfrac{\text{Tr}A(\boldx)}{\vert A(\boldx)\vert^{2}}
 = \dfrac{\Sigma_{1\leq i \leq d}a_{ii}(\boldx)}{\Sigma_{1\leq i,j\leq d}(a_{ij}(\boldx))^{2}}, 
\qquad \forall \boldx \in \Omega.
\end{align}

\begin{align}
\label{nondiv_pde}
\tilde{A}:D^{2}u + \tilde{\boldsymbol{b}}\cdot \nabla u 
+ \tilde{c}u = \tilde{f}  \text{ in } \Omega, \qquad
u = 0 \text{ on } \partial\Omega,
\end{align}

\begin{align}
\label{nondiv_fem}
& \int_{\Omega}(\tilde{A}: \overline{\nabla}_{h}(\nabla u_{h}) + \tilde{\boldsymbol{b}}\cdot \nabla u_{h} 
+ \tilde{c}u_{h} ) v_{h} d\boldx \\
\nonumber 
= & \int_{\Omega} \gamma (A: \overline{\nabla}_{h}(\nabla u_{h}) + \boldsymbol{b}\cdot \nabla u_{h} 
+ cu_{h} ) v_{h} d\boldx 
= \int_{\Omega} \gamma fv_{h} d\boldx = \int_{\Omega} \tilde{f}v_{h} d\boldx, \qquad \forall v_{h} \in V_{h}.
\end{align}

\begin{theorem}
\label{thm_global_A_continuous}
(Global $W^{2,p}$ estimate of linear elliptic PDE in non-divergence form)
Let $A \in [C^{0}(\overline{\Omega})]^{d\times d}$, $\boldsymbol{b}\in 
[L^{\infty}(\Omega)]^{d}$, $c \in L^{\infty}(\Omega)$ with $c \leq 0$, 
 and $\Omega$ is a bounded open Lipschitz polyhedral domain 
in $\mathbb{R}^{d}$ ($d=2,3$).  We further assume $A$ is uniformly elliptic in $\Omega$.
 There is a constant $C_{p}>0$ such that 
\begin{align}
\label{global_w2p_estimate_A_continuous}
\Vert w\Vert_{W^{2,p}(\Omega)} \leq C_{p} \Vert A:D^{2}w  + \boldsymbol{b}\cdot 
\nabla w + c w\Vert_{L^{p}(\Omega)}, 
\qquad \forall w \in W^{2,p}(\Omega) \cap W_{0}^{1,p}(\Omega).
\end{align}
Here $p$ is valid in the range described in (\ref{p_range1}).
\end{theorem}

\begin{lemma}
\label{lemma_conv_compactness} 
(Discrete compactness)
Let $1 < p < + \infty$ and $\Omega$ be a Lipschitz polyhedra in $\mathbb{R}^{d}$ 
($d=2,3$). We assume that there are $0 < M_{1} < +\infty$ and 
$\{w_{h} \in V_{h}\}_{h > 0}$ satisfying 
\begin{align*}
\Vert  w_{h} \Vert_{W_{h}^{2,p}(\Omega)} 
\leq M_{1}, \qquad \forall h > 0.
\end{align*} 
Then there is $w \in W^{2,p}(\Omega) \cap W_{0}^{1,p}(\Omega)$ and $ 0<M_{2}< + \infty$ 
such that 
\begin{align*}
\Vert w_{h} - w\Vert_{W^{1,p}(\Omega)} \rightarrow 0, \qquad 
\overline{\nabla}_{h}(\nabla w_{h}) \rightharpoonup D^{2}w \text{ in } 
[L^{p}(\Omega)]^{d \times d},
\end{align*}
for a subsequence of $h \rightarrow 0$; and 
\begin{align*}
\quad \Vert \overline{\nabla}_{h}(\nabla w) \Vert_{L^{p}(\Omega)} 
+ \Vert \overline{\nabla}_{h} (\nabla w_{h})\Vert_{L^{p}(\Omega)} 
\leq M_{2}, \qquad \forall h> 0.
\end{align*}
\end{lemma}

\begin{theorem}
\label{thm_global_estimate}
(Stability of the FEM (\ref{nondiv_fem})) 
Let $A \in [C^{0}(\overline{\Omega})]^{d\times d}$ uniformly elliptic, $\boldsymbol{b}\in 
[L^{\infty}(\Omega)]^{d}$, $c \in L^{\infty}(\Omega)$ with $c \leq 0$, 
 and $\Omega$ is a bounded open Lipschitz polyhedral domain in $\mathbb{R}^{d}$ ($d=2,3$).  
There is $h_{2} > 0$ which may depend on $p$, such that for all $h \in (0, h_{2})$,
\begin{align}
\label{global_estimate}
\Vert w_{h}\Vert_{W^{1,p}(\Omega)} + \Vert w_{h} \Vert_{W_{h}^{2,p}(\Omega)} 
\leq C \Vert \tilde{A}: \overline{\nabla}_{h}(\nabla w_{h}) +\tilde{\boldsymbol{b}}\cdot\nabla w_{h} 
+ \tilde{c} w_{h}\Vert_{L_{h}^{p}(\Omega)}, \quad \forall w_{h} \in V_{h}.
\end{align}
Here $p$ is valid in the range described in (\ref{p_range2}). 
$\tilde{A}$, $\tilde{\boldsymbol{b}}$ and $\tilde{c}$ are defined in (\ref{nondiv_pde}). 
The norm $\Vert \cdot \Vert_{L_{h}^{p}(\Omega)}$ is defined in (\ref{discrete_zero_order_norm}).
\end{theorem}

\begin{theorem}
\label{thm_A_continuous_complete}
Let $A \in [C^{0}(\overline{\Omega})]^{d\times d}$ uniformly elliptic, 
$\boldsymbol{b}\in [L^{\infty}(\Omega)]^{d}$, $c \in L^{\infty}(\Omega)$ with $c \leq 0$, 
 and $\Omega$ is a bounded open Lipschitz polyhedral domain in $\mathbb{R}^{d}$ ($d=2,3$).  
For any $f \in L^{p}(\Omega)$, 
there is a unique $u \in W^{2,p}(\Omega) \cap W_{0}^{1,p}(\Omega)$ such that 
\begin{subequations}
\label{A_continuous_wellposedness_props}
\begin{align}
\label{A_continuous_wellposedness_prop1}
& A:D^{2} u + \boldsymbol{b}\cdot \nabla u + cu = f \text{ in } \Omega, \\
\label{A_continuous_wellposedness_prop2}
& \Vert u\Vert_{W^{2,p}(\Omega)} \leq C_{p} \Vert f \Vert_{L^{p}(\Omega)}.
\end{align}
\end{subequations}
Here the constant $C_{p}$ is the same as the one in Theorem~\ref{thm_global_A_continuous}. 
$p$ is valid in the range described in (\ref{p_range1}).

Let $u_{h} \in V_{h}$ be the numerical solution of the finite element method 
(\ref{nondiv_fem}). Then we have 
\begin{align}
\label{A_continuous_conv}
\Vert u_{h} -  u \Vert_{W^{1,p}(\Omega)} \rightarrow 0, \qquad 
\overline{\nabla}_{h}(\nabla u_{h}) \rightharpoonup D^{2}u \text{ in } 
[L^{p}(\Omega)]^{d \times d},
\end{align}
as $h \rightarrow 0$. 
And if $h \in (0,h_{2})$ where $h_{2}$ is introduced in Theorem~\ref{thm_global_estimate}, then for any 
$\chi_{h} \in V_{h}$,
\begin{align}
\label{A_continuous_conv_rate}
& \Vert u_{h} - u \Vert_{W^{1,p}(\Omega)}
+ \Vert u_{h} - u \Vert_{W_{h}^{2,p}(\Omega)} \\
\nonumber
\leq C & \big( \Vert u  - \chi_{h} \Vert_{W^{1,p}(\Omega)} + \Vert u - \chi_{h}\Vert_{W_{h}^{2,p}(\Omega)} 
+ \Vert \overline{P}_{h}(D^{2}u) - D^{2}u \Vert_{L^{p}(\Omega)} \big).
\end{align}
Here $p$ is valid in the range described in (\ref{p_range2}). 
$\overline{P}_{h}$ denotes the standard $L^{2}$-orthogonal projection onto $[\overline{V}_{h}]^{d\times d}$.
\end{theorem}

\begin{proof}
We divide the proof into two steps. In step $1$,
We will first prove the estimates (\ref{A_continuous_wellposedness_props}) 
and (\ref{A_continuous_conv},\ref{A_continuous_conv_rate}) for 
 $1< p \leq 2$ if $\Omega$ is convex polyhedra in $\mathbb{R}^{d}$ ($d=2,3$), and 
 $ \frac{4}{3} - \epsilon_{1} <p < \frac{4}{3} 
+ \epsilon_{2}$ if $\Omega$ is a polygon (maybe non-convex) in $\mathbb{R}^{2}$. 
In step $2$, we will prove the estimate (\ref{A_continuous_wellposedness_props}) 
for $1< p \leq \frac{4}{3} - \epsilon_{1}$ and $\Omega$ is a 
polygon (maybe non-convex) in $\mathbb{R}^{2}$. 

Step $1$. We consider $1< p \leq 2$ if $\Omega$ is convex polyhedra in $\mathbb{R}^{d}$ ($d=2,3$), and $ \frac{4}{3} - \epsilon_{1} <p < \frac{4}{3} 
+ \epsilon_{2}$ if $\Omega$ is a polygon (maybe non-convex) in $\mathbb{R}^{2}$.

By the design of the finite element method (\ref{nondiv_fem}) and 
Theorem~\ref{thm_global_estimate}, there is a unique $u_{h} \in V_{h}$ 
to be the numerical solution of (\ref{nondiv_fem}) and 
\begin{align*}
\Vert u_{h} \Vert_{W_{h}^{2,p}(\Omega)} \leq C \Vert f\Vert_{L^{p}(\Omega)},
\end{align*}
if $h$ is small enough. Then by Lemma~\ref{lemma_conv_compactness}, there is 
$u \in W^{2,p}(\Omega) \cap W_{0}^{1,p}(\Omega)$ such that 
\begin{align*}
\Vert u_{h} - u\Vert_{W^{1,p}(\Omega)} \rightarrow 0, \qquad 
\overline{\nabla}_{h}(\nabla u_{h}) \rightharpoonup D^{2}u 
\text{ in } [L^{p}(\Omega)]^{d \times d},
\end{align*}
for a sequence of $h \rightarrow 0$; and there is $0 < M_{2} < +\infty$ such that 
$\Vert \overline{\nabla}(\nabla u_{h})\Vert_{L^{p}(\Omega)} \leq M_{2}$ for any 
$h > 0$.

Now we need to prove (\ref{A_continuous_wellposedness_prop1}). 
We define $p^{\prime} \in (1, + \infty)$ by  $\frac{1}{p} + \frac{1}{p^{\prime}} = 1$. 
We denote by $P_{h}$ the standard $L^{2}$-orthogonal projection onto $V_{h}$.
It is well-known that there is a constant $\tilde{C}_{1}$ independent of $h$, 
\begin{align}
\label{L2_projection_bound3}
\Vert P_{h} v\Vert_{L^{p^{\prime}}(\Omega)} \leq \tilde{C}_{1} \Vert v\Vert_{L^{p^{\prime}}(\Omega)}, 
\qquad \Vert P_{h} v\Vert_{W^{1,p^{\prime}}(\Omega)} \leq \tilde{C}_{1} \Vert v\Vert_{W^{1,p^{\prime}}(\Omega)}.
\end{align}

We choose $v \in C_{0}^{\infty}(\Omega)$ arbitrarily. 
By the design of the finite element method (\ref{nondiv_fem}), We have 
\begin{align*}
& (\tilde{A}:D^{2}u + \tilde{\boldsymbol{b}}\cdot\nabla u 
+ \tilde{c}u, v)_{\Omega} \\
= & (\tilde{A}: (D^{2}u - \overline{\nabla}_{h}(\nabla u_{h})), v)_{\Omega} 
+ (\tilde{A}:\overline{\nabla}_{h}(\nabla u_{h}) 
+ \tilde{\boldsymbol{b}}\cdot\nabla u + \tilde{c}u, v)_{\Omega} \\
= & (\tilde{A}: (D^{2}u - \overline{\nabla}_{h}(\nabla u_{h})), v)_{\Omega} 
+ (\tilde{A}:\overline{\nabla}_{h}(\nabla u_{h}) 
+ \tilde{\boldsymbol{b}}\cdot\nabla u + \tilde{c}u, v - P_{h}v)_{\Omega} \\
& \qquad + (\tilde{A}:\overline{\nabla}_{h}(\nabla u_{h}) 
+ \tilde{\boldsymbol{b}}\cdot\nabla u + \tilde{c}u, P_{h}v)_{\Omega} \\
= & (\tilde{A}: (D^{2}u - \overline{\nabla}_{h}(\nabla u_{h})), v)_{\Omega} 
+ (\tilde{A}:\overline{\nabla}_{h}(\nabla u_{h}) 
+ \tilde{\boldsymbol{b}}\cdot\nabla u + \tilde{c}u, v - P_{h}v)_{\Omega} \\
& \qquad + (\tilde{A}:\overline{\nabla}_{h}(\nabla u_{h}) 
+ \tilde{\boldsymbol{b}}\cdot\nabla u_{h} + \tilde{c}u_{h}, P_{h}v)_{\Omega} 
+ (\tilde{\boldsymbol{b}}\cdot \nabla(u - u_{h}) 
+ \tilde{c}(u - u_{h}), P_{h}v)_{\Omega} \\
= & (\tilde{A}: (D^{2}u - \overline{\nabla}_{h}(\nabla u_{h})), v)_{\Omega} 
+ (\tilde{A}:\overline{\nabla}_{h}(\nabla u_{h}) 
+ \tilde{\boldsymbol{b}}\cdot\nabla u + \tilde{c}u, v - P_{h}v)_{\Omega} \\
& \qquad + (\tilde{f}, v)_{\Omega} + (\tilde{f}, P_{h}v - v)_{\Omega} 
+ (\tilde{\boldsymbol{b}}\cdot \nabla(u - u_{h}) 
+ \tilde{c}(u - u_{h}), P_{h}v)_{\Omega}.
\end{align*}

Since $\overline{\nabla}_{h}(\nabla u_{h}) \rightharpoonup D^{2}u 
\text{ in } [L^{p}(\Omega)]^{d \times d}$, we have 
\begin{align*}
(\tilde{A}:(D^{2}u - \overline{\nabla}_{h}(\nabla u_{h})), v)_{\Omega} 
+ (\tilde{\boldsymbol{b}}\cdot \nabla(u - u_{h}) 
+ \tilde{c}(u - u_{h}), P_{h}v)_{\Omega}\rightarrow 0,
\end{align*}
for a sequence of $h \rightarrow 0$.

By (\ref{L2_projection_bound3}) and the fact that $\Vert \overline{\nabla}(\nabla u_{h})\Vert_{L^{p}(\Omega)} 
\leq M_{2}$, we have 
\begin{align*}
& \vert (\tilde{A}:\overline{\nabla}_{h}(\nabla u_{h})
+ \tilde{\boldsymbol{b}}\cdot \nabla u + \tilde{c}u, v - P_{h}v)_{\Omega} \vert \\
\leq & C \left(\Vert \overline{\nabla}_{h}(\nabla u_{h})\Vert_{L^{p}(\Omega)} 
+ \Vert u\Vert_{W^{1,p}(\Omega)}\right)\Vert v - P_{h} 
v\Vert_{L^{p^{\prime}}(\Omega)} \rightarrow 0,
\end{align*}
as $h \rightarrow 0$. 

Similarly, we have 
\begin{align*}
\vert (\tilde{f}, P_{h}v - v)_{\Omega} \vert \rightarrow 0,
\end{align*}
as $h \rightarrow 0$.

So we have 
\begin{align*}
(\tilde{A}:D^{2}u + \tilde{\boldsymbol{b}}\cdot\nabla u 
+ \tilde{c}u, v)_{\Omega} = (\tilde{f}, v)_{\Omega}. 
\end{align*}

Since $v \in C_{0}^{\infty}(\Omega)$ is chosen arbitrarily, 
$\tilde{A}:D^{2}+\tilde{\boldsymbol{b}}\cdot\nabla u 
+ \tilde{c}u = \tilde{f}$ in $\Omega$ almost everywhere. 
By the definition of the function $\gamma$ (see (\ref{def_gamma})) and 
the fact that $A$ is uniformly elliptic and uniformly bounded in $\Omega$,  
$\gamma$ is uniformly bounded from below and above by positive constants. 
So we have $A:D^{2}u +\boldsymbol{b}\cdot\nabla u + cu = f$ in $\Omega$ 
almost everywhere. 
So (\ref{A_continuous_wellposedness_prop1}) 
is proven. Theorem~\ref{thm_global_A_continuous} implies the uniqueness of 
$u \in W^{2,p}(\Omega) \cap W_{0}^{1,p}(\Omega)$, which results in 
\begin{align*}
\Vert u_{h} - u\Vert_{W^{1,p}(\Omega)} \rightarrow 0, \qquad 
\overline{\nabla}_{h}(\nabla u_{h}) \rightharpoonup D^{2}u 
\text{ in } [L^{p}(\Omega)]^{d \times d},
\end{align*}
as $h \rightarrow 0$. 

The proof of (\ref{A_continuous_conv_rate}) is straightforward due to 
Theorem~\ref{thm_global_estimate} and the construction of the finite element method (\ref{nondiv_fem}).
Thus the proof in Step $1$ is complete.

Step $2$. We consider $1< p \leq \frac{4}{3} - \epsilon_{1}$ if $\Omega$ is a 
polygon (maybe non-convex) in $\mathbb{R}^{2}$. 

Let $\{{f_{i}}\}_{i \geq 1} \subset L^{\frac{4}{3}}(\Omega)$ such that 
$\Vert f_{i} -f \Vert_{L^{p}(\Omega)} \rightarrow 0$ as $i \rightarrow +\infty$. 
For any $i \geq 1$, we denote by $u_{i} \in W^{2,\frac{4}{3}}(\Omega) \cap 
W_{0}^{1, \frac{4}{3}}(\Omega)$ satisfying 
\begin{align*}
A:D^{2} u_{i} + \boldsymbol{b}\cdot \nabla u_{i} + c u_{i} = f_{i} 
\text{ in } \Omega \text{ almost everywhere}.
\end{align*}

Since $\Vert f_{i} -f \Vert_{L^{p}(\Omega)} \rightarrow 0$ as $i \rightarrow +\infty$, 
$\{ \Vert f_{i} \Vert_{L^{p}(\Omega)} \}_{i \geq 1}$ have a uniform upper bound. 
So $\{ \Vert u_{i}\Vert_{W^{2,p}(\Omega)} \}_{i \geq 1}$ have a uniform upper bound 
as well. So there is $u \in W^{2,p}(\Omega) \cap W_{0}^{1,p}(\Omega)$ such that 
\begin{align*}
u_{i} \rightharpoonup u \text{ in } W^{2,p}(\Omega),
\end{align*}
for a subsequence of $i \rightarrow + \infty$.

We choose $v \in C_{0}^{\infty}(\Omega)$ arbitrarily. Then 
\begin{align*}
& (A:D^{2} u + \boldsymbol{b}\cdot \nabla u + cu, v)_{\Omega} \\
= & (A:D^{2}(u - u_{i}) + \boldsymbol{b}\cdot \nabla (u - u_{i}) 
+ c(u - u_{i}), v)_{\Omega} 
+ (A:D^{2}u_{i} + \boldsymbol{b}\cdot\nabla u_{i} + c u_{i}, v)_{\Omega} \\
= & (A:D^{2}(u - u_{i}) + \boldsymbol{b}\cdot \nabla (u - u_{i}) 
+ c(u - u_{i}), v)_{\Omega} + (f_{i} - f, v)_{\Omega}
+ (f, v)_{\Omega}. 
\end{align*}

Since $u_{i} \rightharpoonup u \text{ in } W^{2,p}(\Omega)$ for 
a subsequence of $i \rightarrow + \infty$ and 
$\Vert f_{i} - f \Vert_{L^{p}(\Omega)} \rightarrow 0$ as 
$i \rightarrow + \infty$, we have 
\begin{align*}
(A:D^{2} u + \boldsymbol{b}\cdot \nabla u 
+ c u, v)_{\Omega} = (f, v)_{\Omega}.
\end{align*}

Since $v \in C_{0}^{\infty}(\Omega)$ is chosen arbitrarily, 
$A:D^{2}u +  \boldsymbol{b}\cdot \nabla u + c u= f$ in $\Omega$ almost everywhere. 
By Theorem~\ref{thm_global_A_continuous}, 
$u \in W^{2,p}(\Omega)\cap W_{0}^{1,p}(\Omega)$ is the unique solution of 
(\ref{nondiv_pde_original}). Thus the proof of Step $2$ is complete. 
\end{proof}

\end{document}
